% ==================== 2.8. CONCLUSIONES DEL ESTADO DEL ARTE ====================
\section{Conclusiones}\label{sec:conclusiones_soa}

La revisión bibliográfica realizada evidencia de forma concluyente el notable avance de la Inteligencia Artificial aplicada al análisis de la voz como herramienta para la detección temprana de enfermedades neurodegenerativas. El análisis de los doce estudios seleccionados demuestra que la bioseñal vocal contiene información acústica, prosódica y articulatoria capaz de reflejar, con alta sensibilidad, las alteraciones motoras y cognitivas inherentes a patologías como la Enfermedad de Parkinson, el Alzheimer y la Esclerosis Lateral Amiotrófica (ELA).

Los resultados reportados en la literatura, con precisiones diagnósticas que frecuentemente superan el 95\%, ratifican que la voz se consolida como un biomarcador digital no invasivo, económico y de alta accesibilidad. A nivel algorítmico, los métodos de aprendizaje automático clásico han demostrado una notable eficacia en corpus reducidos apoyados en una selección rigurosa de descriptores. Paralelamente, las redes neuronales profundas y los modelos fundacionales han inaugurado un nuevo paradigma, ofreciendo una capacidad superior de generalización y la habilidad de inferir patrones latentes complejos directamente desde la señal acústica o sus representaciones espectrales.

No obstante, la disciplina se encuentra aún en una fase de maduración tecnológica. Subsisten limitaciones estructurales significativas que deben ser solventadas, tales como la carencia de bases de datos públicas, masivas y estandarizadas, la escasa validación clínica sobre cohortes externas, la alta varianza en los protocolos de captura y la opacidad algorítmica de las redes profundas. Superar estos desafíos metodológicos constituye el hito fundamental para garantizar la fiabilidad, equidad y aplicabilidad real de los sistemas desarrollados.

Las tendencias tecnológicas actuales trazan una hoja de ruta clara orientada hacia la adopción de modelos auto-supervisados, la integración multimodal (fusión de voz, texto y datos clínicos), el desarrollo de técnicas de Inteligencia Artificial Explicable (XAI) y la creación de repositorios en abierto que garanticen la reproducibilidad científica. En conjunto, estas directrices configuran un horizonte investigador excepcionalmente prometedor. La convergencia interdisciplinar entre la ingeniería de datos, la lingüística computacional y la práctica médica permitirá que los sistemas predictivos trasciendan el entorno experimental de laboratorio, integrándose como verdaderas herramientas de apoyo a la decisión clínica para la monitorización y el diagnóstico precoz.